% GENERATED FILE. Edit canonical lessons, then run npm run generate:books.
% canonical-source-hash: fnv1a64:097cdf773fa3008c
% canonical-lessons: JA-C206-norimono, JA-C206-koutsuu, JA-C206-juutai, JA-C206-joukyaku, JA-C206-yoyaku

\chapter{Traffic, Passengers and Bookings}
\label{ch:ja-traffic-passengers-and-bookings}
\hlchaptermodality{206}

\begin{chapteropening}
\textbf{By the end of this chapter:} \emph{I can talk about traffic, passengers and bookings.}

Complete the last lesson of chapter 206: I can talk about traffic, passengers and bookings.
\end{chapteropening}

\section[\texorpdfstring{\ja{のりもの} (norimono)}{のりもの (norimono)}]{\ja{のりもの} (norimono) — a vehicle; a means of transport}

\label{lesson:JA-C206-norimono}



\begin{quote}
Before the new one: say the Japanese for one way, then the Japanese for a round trip.
\end{quote}

\subsection*{You'll want to know: \ja{のりもの}}
\textbf{\ja{のりもの}} — \emph{norimono} — \textquotedblleft{}a vehicle; a means of transport\textquotedblright{}.

\begin{quote}
\textbf{\ja{でんしゃや} \ja{ふねは} \ja{のりものです。}} — \emph{densha ya fune wa norimono desu} — \textquotedblleft{}Trains and boats are vehicles.\textquotedblright{}
\end{quote}

\begin{cousinweb}
\textbf{\ja{のりもの}} (\emph{norimono}) is native Japanese (wago), built from \emph{nori}, \textquotedblleft{}riding\textquotedblright{}, and \emph{mono}, \textquotedblleft{}thing\textquotedblright{}, both inherited rather than borrowed (JMdict).
\end{cousinweb}

\begin{grammarlens}[title={What you've built}]
One more word for everyday life at A2.
\end{grammarlens}

\subsection*{Your turn}
\begin{itemize}
\raggedright
  \item \emph{Say it:} \emph{norimono}
  \item \emph{Say it:} \emph{norimono}, once more
  \item \emph{Say it:} \emph{densha ya fune wa norimono desu} — Trains and boats are vehicles.
  \item \emph{Recall:} say the Japanese for one way, then the Japanese for a round trip, then say \emph{norimono} again
\end{itemize}

\begin{tcolorbox}[breakable,colback=teal!4,colframe=teal!35!black,title={Before you move on}]
What does \ja{のりもの} mean? (\textquotedblleft{}A vehicle; a means of transport\textquotedblright{}.) Say it once more.
\end{tcolorbox}
\section[\texorpdfstring{\ja{こうつう} (kōtsū)}{こうつう (kōtsū)}]{\ja{こうつう} (kōtsū) — traffic; transport}

\label{lesson:JA-C206-koutsuu}



\begin{quote}
Before the new one: say the Japanese for a round trip, then the Japanese for a vehicle.
\end{quote}

\subsection*{You'll want to know: \ja{こうつう}}
\textbf{\ja{こうつう}} — \emph{kōtsū} — \textquotedblleft{}traffic; transport\textquotedblright{}.

\begin{quote}
\textbf{\ja{この} \ja{みちは} \ja{こうつうが} \ja{おおいです。}} — \emph{kono michi wa kōtsū ga ōi desu} — \textquotedblleft{}This road has a lot of traffic.\textquotedblright{}
\end{quote}

\begin{cousinweb}
\textbf{\ja{こうつう}} (\emph{kōtsū}) is Sino-Japanese (kango): two Chinese characters read with their on'yomi, readings borrowed from Chinese. They are \emph{kō}, \textquotedblleft{}to cross, mix\textquotedblright{} (Middle Chinese \emph{kaew}) and \emph{tsū}, \textquotedblleft{}to pass through\textquotedblright{} (\emph{thuwng}) (KANJIDIC2; Middle Chinese in Baxter's transcription).
\end{cousinweb}

\begin{grammarlens}[title={What you've built}]
One more word for everyday life at A2.
\end{grammarlens}

\subsection*{Your turn}
\begin{itemize}
\raggedright
  \item \emph{Say it:} \emph{kōtsū}
  \item \emph{Say it:} \emph{kōtsū}, once more
  \item \emph{Say it:} \emph{kono michi wa kōtsū ga ōi desu} — This road has a lot of traffic.
  \item \emph{Recall:} say the Japanese for a round trip, then the Japanese for a vehicle, then say \emph{kōtsū} again
\end{itemize}

\begin{tcolorbox}[breakable,colback=teal!4,colframe=teal!35!black,title={Before you move on}]
What does \ja{こうつう} mean? (\textquotedblleft{}Traffic; transport\textquotedblright{}.) Say it once more.
\end{tcolorbox}
\section[\texorpdfstring{\ja{じゅうたい} (jūtai)}{じゅうたい (jūtai)}]{\ja{じゅうたい} (jūtai) — a traffic jam; congestion}

\label{lesson:JA-C206-juutai}



\begin{quote}
Before the new one: say the Japanese for a vehicle, then the Japanese for traffic.
\end{quote}

\subsection*{You'll want to know: \ja{じゅうたい}}
\textbf{\ja{じゅうたい}} — \emph{jūtai} — \textquotedblleft{}a traffic jam; congestion\textquotedblright{}.

\begin{quote}
\textbf{\ja{あさは} \ja{どうろが} \ja{じゅうたいです。}} — \emph{asa wa dōro ga jūtai desu} — \textquotedblleft{}The roads are jammed in the morning.\textquotedblright{}
\end{quote}

\begin{cousinweb}
\textbf{\ja{じゅうたい}} (\emph{jūtai}) is Sino-Japanese (kango): two Chinese characters read with their on'yomi, readings borrowed from Chinese. They are \emph{jū}, \textquotedblleft{}astringent, reluctant\textquotedblright{} (Middle Chinese \emph{srip}) and \emph{tai}, \textquotedblleft{}to stagnate\textquotedblright{} (\emph{drjejH}) (KANJIDIC2; Middle Chinese in Baxter's transcription).
\end{cousinweb}

\begin{grammarlens}[title={What you've built}]
One more word for everyday life at A2.
\end{grammarlens}

\subsection*{Your turn}
\begin{itemize}
\raggedright
  \item \emph{Say it:} \emph{jūtai}
  \item \emph{Say it:} \emph{jūtai}, once more
  \item \emph{Say it:} \emph{asa wa dōro ga jūtai desu} — The roads are jammed in the morning.
  \item \emph{Recall:} say the Japanese for a vehicle, then the Japanese for traffic, then say \emph{jūtai} again
\end{itemize}

\begin{tcolorbox}[breakable,colback=teal!4,colframe=teal!35!black,title={Before you move on}]
What does \ja{じゅうたい} mean? (\textquotedblleft{}A traffic jam; congestion\textquotedblright{}.) Say it once more.
\end{tcolorbox}
\section[\texorpdfstring{\ja{じょうきゃく} (jōkyaku)}{じょうきゃく (jōkyaku)}]{\ja{じょうきゃく} (jōkyaku) — a passenger}

\label{lesson:JA-C206-joukyaku}



\begin{quote}
Before the new one: say the Japanese for traffic, then the Japanese for a traffic jam.
\end{quote}

\subsection*{You'll want to know: \ja{じょうきゃく}}
\textbf{\ja{じょうきゃく}} — \emph{jōkyaku} — \textquotedblleft{}a passenger\textquotedblright{}.

\begin{quote}
\textbf{\ja{ふねの} \ja{じょうきゃくは} \ja{すくないです。}} — \emph{fune no jōkyaku wa sukunai desu} — \textquotedblleft{}The ship has few passengers.\textquotedblright{}
\end{quote}

\begin{cousinweb}
\textbf{\ja{じょうきゃく}} (\emph{jōkyaku}) is Sino-Japanese (kango): two Chinese characters read with their on'yomi, readings borrowed from Chinese. They are \emph{jō}, \textquotedblleft{}to ride\textquotedblright{} (Middle Chinese \emph{zying}) and \emph{kyaku}, \textquotedblleft{}guest\textquotedblright{} (\emph{khaek}) (KANJIDIC2; Middle Chinese in Baxter's transcription).
\end{cousinweb}

\begin{grammarlens}[title={What you've built}]
One more word for everyday life at A2.
\end{grammarlens}

\subsection*{Your turn}
\begin{itemize}
\raggedright
  \item \emph{Say it:} \emph{jōkyaku}
  \item \emph{Say it:} \emph{jōkyaku}, once more
  \item \emph{Say it:} \emph{fune no jōkyaku wa sukunai desu} — The ship has few passengers.
  \item \emph{Recall:} say the Japanese for traffic, then the Japanese for a traffic jam, then say \emph{jōkyaku} again
\end{itemize}

\begin{tcolorbox}[breakable,colback=teal!4,colframe=teal!35!black,title={Before you move on}]
What does \ja{じょうきゃく} mean? (\textquotedblleft{}A passenger\textquotedblright{}.) Say it once more.
\end{tcolorbox}
\section[\texorpdfstring{\ja{よやく} (yoyaku)}{よやく (yoyaku)}]{\ja{よやく} (yoyaku) — a booking; a reservation}

\label{lesson:JA-C206-yoyaku}



\begin{quote}
Before the new one: say the Japanese for a traffic jam, then the Japanese for a passenger.
\end{quote}

\subsection*{You'll want to know: \ja{よやく}}
\textbf{\ja{よやく}} — \emph{yoyaku} — \textquotedblleft{}a booking; a reservation\textquotedblright{}.

\begin{quote}
\textbf{\ja{よやくは} \ja{ひつようですか。}} — \emph{yoyaku wa hitsuyō desu ka} — \textquotedblleft{}Is a booking necessary?\textquotedblright{}
\end{quote}

\begin{cousinweb}
\textbf{\ja{よやく}} (\emph{yoyaku}) is Sino-Japanese (kango): two Chinese characters read with their on'yomi, readings borrowed from Chinese. They are \emph{yo}, \textquotedblleft{}beforehand\textquotedblright{} (Middle Chinese \emph{yoH}) and \emph{yaku}, \textquotedblleft{}to bind, promise\textquotedblright{} (\emph{'jak}) (KANJIDIC2; Middle Chinese in Baxter's transcription).
\end{cousinweb}

\begin{grammarlens}[title={What you've built}]
One more word for everyday life at A2.
\end{grammarlens}

\subsection*{Your turn}
\begin{itemize}
\raggedright
  \item \emph{Say it:} \emph{yoyaku}
  \item \emph{Say it:} \emph{yoyaku}, once more
  \item \emph{Say it:} \emph{yoyaku wa hitsuyō desu ka} — Is a booking necessary?
  \item \emph{Recall:} say the Japanese for a traffic jam, then the Japanese for a passenger, then say \emph{yoyaku} again
\end{itemize}

\begin{tcolorbox}[breakable,colback=teal!4,colframe=teal!35!black,title={Before you move on}]
What does \ja{よやく} mean? (\textquotedblleft{}A booking; a reservation\textquotedblright{}.) Say it once more.
\end{tcolorbox}
